\documentclass[11pt,a4paper]{article}
\usepackage[T1]{fontenc}
\usepackage{lmodern}
\usepackage[margin=2.5cm]{geometry}
\usepackage{graphicx}
\usepackage{authblk}
\usepackage[numbers,sort&compress]{natbib}
\usepackage[hidelinks]{hyperref}

\title{ToolName: what the software does, in one line}
\author[1]{First Author}
\author[2]{Second Author}
\affil[1]{Department of Bioinformatics, University, City, Country}
\affil[2]{Institute, City, Country}
\date{}

\begin{document}

\maketitle

\begin{abstract}
\noindent\textbf{Summary:} What the tool does and why it is needed, in two or three sentences.\\
\textbf{Availability and implementation:} Where it can be downloaded, its licence and the platforms it runs on.\\
\textbf{Contact:} first@example.org
\end{abstract}

\section{Introduction}
State the need the tool meets and the existing tools~\cite{altschul1990}.

\section{Implementation}
Describe the design and the algorithms.

\section{Results}
Show what the tool does on a typical data set (Figure~\ref{fig:overview})~\cite{jumper2021}.

\begin{figure}[htbp]
  \centering
  \includegraphics[width=0.6\linewidth]{example-image}
  \caption{Overview of the method. Upload an image and use its file name here.}
  \label{fig:overview}
\end{figure}

\section{Conclusion}
Summarise what the tool offers.

\section*{Funding}
Name the funders.

\bibliographystyle{unsrtnat}
\bibliography{references}

\end{document}
